\prefacesection{Abstract}
A popular logic-based combinatorial problem known as Sudoku is computationally expensive because of its interdependent constraints. From traditional solvers to heuristic approaches, although effective for smaller puzzles, face difficulties with more complex puzzles. This paper presents SudoSLVRR, a parallel multi-colony ant optimization framework as a Sudoku solver that addresses scalability limitations of the existing solvers. This framework combines constraint propagation and multi-colony ant optimization with dynamic collaborative mechanism (DCM-ACO) inside the multithreaded environment. Heterogeneous ACS and MMAS colonies are used to balance exploitation and exploration. This multi-colony ACO is executed inside each thread where threads exchange their iteration-best solution and best-so-far solution in ring and random communication topologies, respectively. The communication between threads prevents premature stagnation and promotes diversity. Experimental results on $9 \times 9$, $16 \times 16$, and $25 \times 25$ Sudoku puzzles show that the SudoSLVRR achieved a higher success rate (100\%) and reduced runtime and iteration count by up to 93\% and 97\% respectively, compared to CP-ACS and CP-DCM-ACO baselines. This result shows that the communication between threads and the collaboration between colonies inside each thread allows the proposed solver to increase its solution diversity which prevents global stagnation, leading to better performance in constraint-driven problems. 


%Sudoku is a logic-based, constraint-satisfaction combinatorial problem that makes it computationally expensive for any traditional solver. Most of the heuristic approaches, while effective for small instances, have problems with scalability, convergence speed, and maintaining solution diversity in larger grids. This article presents a Collaborative Multithreaded Multi-Colony Ant  Optimization for overcoming these challenges. The model embeds Constraint Propagation for the early elimination of infeasible values, a Dynamic Collaborative Multi-Colony structure for adaptive pheromone sharing, and a Randomly Matched communication strategy to enhance cooperation between threads and avoid premature convergence. Experimental results on $9 \times 9$, $16 \times 16$, and $25 \times 25$ Sudoku puzzles shows that the proposed multithread multi-colony ant solver was able to achieve higher success rate and less execution time when compared to existing solvers.

% \noindent \textbf{Keywords:} Constraint Satisfaction Problem, Swarm Intelligence, Parallelization, Ant Colony System, Dynamic Collaborative Mechanism